\documentclass[12pt, oneside]{article}
\usepackage[letterpaper, margin=1in]{geometry}
\usepackage[english]{babel}
\usepackage[utf8]{inputenc}
\usepackage{amsmath}
\usepackage{amsfonts}
\usepackage{amssymb}
\usepackage{tikz}

\usepackage{fancyhdr}
\pagestyle{fancy}
\fancyhf{}
\rhead{\thepage \\Name: \hspace{1.5in}.\\}
\lhead{BECA / Dr. Huson / 11.1 IB Math SL \\* 11 June 2018 \\* \textbf{Classwork: Regents review}}

\vspace{1cm}
\renewcommand{\headrulewidth}{0pt}

\title{Problem set template}
\author{Chris Huson}
\date{June 2018}

\begin{document}
\begin{enumerate}
\item Given the function $f$ represented in the table below.\\*
\begin{center}
\begin{tabular}{|c|c|}
    \hline 
    $\boldsymbol{x}$ & $\boldsymbol{f(x)}$\\ 
    \hline 
    -2 & 5 \\ 
    \hline 
    -1 & 0 \\ 
    \hline 
    0 & 1 \\ 
    \hline 
    1 & -2 \\ 
    \hline 
    2 & 2 \\ 
    \hline 
\end{tabular}\\*[5pt]
\end{center}
\begin{enumerate}
    \item What is the $x$-intercept?\\[25pt]
    \item For what $x$ is $f(x)$ minimum?\\[25pt]
    \item What is the $y$-intercept of the function?\\[15pt]
\end{enumerate}

\newpage

\item A \$4,250 investment earning a continuous interest rate of 3.25\% over 3 years would be worth how much?\\[2in]


\item Explain how $\displaystyle \left(3^{\frac{1}{3}} \right)^2$ can be written as the equivalent radical expression $\sqrt[3]9$. \\[3in] %Alg2 Regents Aug2016

\item Given $i$ is the imaginary unit, simplify $(2x-yi)^2$ to the form $a+bi$. \\[1.8in]%Alg2 Regents Jun2016

\newpage

\item Consider the function $h(x) = 2\sin(x) + 1$.\\*
What is the minimum value of $h(x)$ for the domain $[0,2\pi]$? Justify your answer by sketching a graph. \\[2in]%Alg2 Regents Jan2018


\item Sketch a graph of a cubic polynomial with the following characteristics: 
\begin{itemize}
\item three zeros: the origin, one positive, one negative
\item as $x \rightarrow + \infty$, $f(x) \rightarrow - \infty$
\item as $x \rightarrow - \infty$, $f(x) \rightarrow + \infty$
\end{itemize}
\begin{center}
    \begin{tikzpicture}[scale=2/4]
    \draw[thick,<->] (-7.5,0) -- (7.5,0) node[anchor=north west] {\textbf{x}};
    \draw[thick,<->] (0,-7.5) -- (0,7.5) node[anchor=south east] {\textbf{y}};
    \end{tikzpicture}
\end{center} %Alg2 Regents Jun2016 MC
What is the sign of the function's leading coefficient?




\newpage
\item On the axes below, graph one cycle of a cosine function with amplitude 3, period $\pi$, midline $y=-1$, and passing through the point $(0,2)$.
\begin{center}
    \begin{tikzpicture}[scale=2.54/4]
    \draw[step=1cm,gray,very thin] (-10,-10) grid (10,10);
    \draw[thick,<->] (-10.5,0) -- (10.5,0) node[anchor=north west] {\textbf{x}};
    \draw[thick,<->] (0,-10.5) -- (0,10.5) node[anchor=south east] {\textbf{y}};
    \end{tikzpicture} 
\end{center} %Jun2016 Regents FR
What is the minimum of the function?

\newpage

\item Graph $y= \log_2(x+2)-4$ on the set of axes below. Use an appropriate scale to include both intercepts.
\begin{center}
    \begin{tikzpicture}[scale=2.54/4]
    \draw[step=1cm,gray,very thin] (-10,-10) grid (10,10);
    \draw[thick,<->] (-10.5,0) -- (10.5,0) node[anchor=north west] {\textbf{x}};
    \draw[thick,<->] (0,-10.5) -- (0,10.5) node[anchor=south east] {\textbf{y}};
    \end{tikzpicture}
\end{center}

Describe the behavior of the given function as $x$ approaches $-2$ and as $x$ approaches positive infinity. %Alg2 Regents Jun2017 


\end{enumerate}
\end{document}


\item The value of a certain small passenger car based on its use in years is modeled by $V(t) =28482.698(0.684)^t$, where $V(t)$ is the value in dollars and $t$ is the time in years. Zach had to take out a loan to purchase the small passenger car. The function $Z(t)=22151.327(0.778)^t$, where $Z(t)$ is measured in dollars, and $t$ is the time in years, models the unpaid amount of Zach’s loan over time.\\*[10pt]
Graph $V(t)$ and $Z(t)$ over the interval $0 \leq t \leq 5$, on the set of axes below.
\begin{center}
    \begin{tikzpicture}
    \draw[step=0.25in,gray,very thin] (0,0) grid (12.7,12.7);
    \draw[thick,->] (0,0) -- (13,0);% node[anchor=north west] {x};
    \draw[thick,->] (0,0) -- (0,13);% node[anchor=south east] {y};
    \end{tikzpicture}
\end{center}
State when $V(t)=Z(t)$, to the \emph{nearest hundredth}, and interpret its meaning in the context of the problem.\\*[10pt]
 %Alg2 Regents Aug2017
